\documentclass[12pt,a4paper]{report}

\usepackage[utf8]{inputenc} % pentru suport diacritice
% \usepackage[romanian]{babel} % setări pentru limba română
\renewcommand\familydefault{\sfdefault} % sans serif

\usepackage[margin=2.54cm]{geometry} % dimensiuni pagină și margini
\usepackage{graphicx} % support the \includegraphics command and options

% formatting sections and subsections
\usepackage{textcase}
\usepackage[titletoc, title]{appendix}
\usepackage{titlesec}
\titleformat{\chapter}{\large\bfseries\MakeUppercase}{\thechapter}{2ex}{}[\vspace*{-1.5cm}]
\titleformat*{\section}{\large\bfseries}
\titleformat*{\subsection}{\large\bfseries}
\titleformat*{\subsubsection}{\large\bfseries}

\usepackage{chngcntr}
\counterwithout{figure}{chapter} % no chapter number in figure labels
\counterwithout{table}{chapter} % no chapter number in table labels
\counterwithout{equation}{chapter} % no chapter number in equation labels

\usepackage{booktabs} % for much better looking tables
\usepackage{url} % Useful for inserting web links nicely
\usepackage[bookmarks,unicode,hidelinks]{hyperref}

\usepackage{array} % for better arrays (eg matrices) in maths
\usepackage{paralist} % very flexible & customisable lists (eg. enumerate/itemize, etc.)
\usepackage{verbatim} % adds environment for commenting out blocks of text & for better verbatim
\usepackage{subfig} % make it possible to include more than one captioned figure/table in a single float
\usepackage{enumitem}
\setlist{noitemsep}

% packages used for citations
\usepackage{apacite}
\usepackage{natbib}

\usepackage{xcolor}
\usepackage{listings}
\lstset{language=Python,basicstyle={\bfseries \color{black}},keywordstyle={\bfseries \color{blue}}}

%%% HEADERS & FOOTERS
\usepackage{fancyhdr}
\pagestyle{empty}
\renewcommand{\headrulewidth}{0pt}
\renewcommand{\footrulewidth}{0pt}
\lhead{}\chead{}\rhead{}
\lfoot{}\cfoot{\thepage}\rfoot{}

\newcommand{\HeaderLineSpace}{-0.25cm}

\newcommand{\UniTextEN}{National University of Science and Technology\\ POLITEHNICA Bucharest\\
Faculty of Automatic Control and Computers\\
Computer Science and Engineering Department\\}
\newcommand{\DiplomaEN}{DIPLOMA PROJECT}
\newcommand{\AdvisorEN}{Thesis advisors:}
\newcommand{\BucEN}{BUCHAREST}

\newcommand{\frontPage}[6]{
\begin{titlepage}
\begin{center}
{\Large #1}  % header (university, faculty, department)
\vspace{50pt}
\begin{tabular}{ccc}
\includegraphics[height=3cm]{pics/UNST Poli.png} & 
\includegraphics[height=3cm]{pics/A&C.png} & 
\includegraphics[height=3cm]{pics/UPB NLP.png}
\end{tabular}

\vspace{105pt}
{\Huge #2}\\                             % diploma project text
\vspace{40pt}
{\Large #3}\\ \vspace{0pt}               % project title
\vspace{40pt}
{\LARGE \Name}\\                         % student name
\end{center}
\vspace{60pt}
\begin{tabular*}{\textwidth}{@{\extracolsep{\fill}}p{6cm}r}
&{\large\textbf{#5}}\vspace{10pt}\\      % scientific advisor
&{\large \AdvisorA}\\                    % advisor name A
&{\large \AdvisorB}                      % advisor name B
\end{tabular*}
\vspace{20pt}
\begin{center}
{\large\textbf{#6}}\\
\vspace{0pt}
{\normalsize \Year}
\end{center}
\end{titlepage}
}

\newcommand{\frontPageEN}{\frontPage{\UniTextEN}{\DiplomaEN}{\ProjectTitleEN}{\ProjectSubtitleEN}{\AdvisorEN}{\BucEN}}

\linespread{1.15}
\setlength\parindent{0pt}
\setlength\parskip{.28cm}

%% Abstract macro
\newcommand{\AbstractPage}{
\begin{titlepage}

\thispagestyle{fancy}\setcounter{page}{5}
\clearpage\phantomsection\addcontentsline{toc}{chapter}{Abstract}

\textbf{\large ABSTRACT}\par
\AbstractEN \vfill
\textbf{\large REZUMAT}\par
\AbstractRO\par\vfill
\end{titlepage}
}


%%%%%%%%%%%%%%%%%%%%%%%%%%%%%%%%%%%%%%%%%%%%%%%%%%
%%
%%          End of template definitions
%%
%%%%%%%%%%%%%%%%%%%%%%%%%%%%%%%%%%%%%%%%%%%%%%%%%%

\newcommand{\ProjectTitleEN}{Awesome Title}
\newcommand{\Name}{Your name}
\newcommand{\AdvisorA}{Prof./Conf./S.l./As. dr./drd. ing. X}
\newcommand{\AdvisorB}{Prof./Conf./S.l./As. dr./drd. ing. Y}
\newcommand{\Year}{202x}

% Setări document
\title{Proiect de diplomă}
\author{\Name}
\date{\Year}

%%
%%   Câmpurile aferente rezumatului
%%
\newcommand{\AbstractEN}{
The abstract should contain the following information (1/2 sentences per item):
\begin{itemize}
    \item Problem statement / context
    \item Aim of the thesis/research objective(s)
    \item Brief method details
    \item Main results
    \item Concluding remarks
\end{itemize}
}

\newcommand{\AbstractRO}{Translation to be made at the end

\vspace{3\baselineskip}
\textbf{Keywords}: max 5, e.g., Language Models, Natural Language Processing
}


\begin{document}
\pagenumbering{roman}
\frontPageEN
\begingroup
\linespread{1}
\tableofcontents
\endgroup

% Keep the list of figures and the list of tables on the same page.
\pagestyle{fancy}
\listoffigures\addcontentsline{toc}{chapter}{List of figures}
\begingroup
\let\clearpage\relax
\listoftables\addcontentsline{toc}{chapter}{List of tables}
\endgroup

% Include the abstract page. See this definition if page numbers are wrong.
\AbstractPage

% poate fi comentata sau stearsa
% \ThanksPage

% Start the normal page numbering.
\pagenumbering{arabic}

% Textul licentei incepe de aici
\chapter{Introduction}\pagestyle{fancy}

Google Scholar and Mendeley are your friends for generating references. You can also export them directly from DBLP and other websites.

Please use $citet$ - e.g., \citet{NIPS2017_3f5ee243} - for contextual references to the authors and $citep$ - e.g., \citep{HochSchm97} - in usual citations. \textbf{Never} include paper titles, only formulations that relate to the authors and are automatically generated (e.g., The study by ..., or X (year) have ...). Examples of BibTex entries are included in the "bibliography.bib" file.

Also, Grammarly is extremely useful for grammatical corrections. Please use the Chrome plugin integrated with Overleaf.

\section{Context}

\section{Problem Statement}

\section{Objectives}

\section{Thesis Structure}


%---------------------------
\chapter{State of the art}



%---------------------------
\chapter{Method}

\section{Corpus}

\subsection{Descriptives}

\section{(Neural) Architecture}

\section{Performance Metrics}


%---------------------------
\chapter{Results}


%---------------------------
\chapter{Discussion}

\section{Performance Comparison}

\section{Limitations}


%---------------------------
\chapter{Conclusions and Future Work}


% Bibliography / References.
\renewcommand\bibname{References}
\clearpage\phantomsection\addcontentsline{toc}{chapter}{\bibname}
\bibliographystyle{apacite}
\bibliography{bibliography}

\end{document}
